% Generated draft: requires mathematical and source review.
\documentclass[11pt]{article}
\usepackage[margin=25mm]{geometry}
\usepackage{fontspec}
\setmainfont{Latin Modern Roman}
\usepackage{amsmath,amssymb,mathtools,unicode-math}
\setmathfont{Latin Modern Math}
\usepackage{enumitem,needspace,xurl,hyperref}
\setlength{\parindent}{0pt}
\setlength{\parskip}{5pt}
\setlength{\emergencystretch}{3em}
\newcommand{\N}{\mathbb{N}}
\newcommand{\Z}{\mathbb{Z}}
\newcommand{\Q}{\mathbb{Q}}
\newcommand{\R}{\mathbb{R}}
\newcommand{\C}{\mathbb{C}}
\newcommand{\F}{\mathbb{F}}
\newcommand{\A}{\mathbb{A}}
\newcommand{\PP}{\mathbb{P}}
\newcommand{\E}{\mathbb{E}}
\newcommand{\eps}{\varepsilon}
\newcommand{\dd}{\,\mathrm{d}}
\begin{document}
\section*{UnsolvedMath ID: 30001390}
\textbf{Source problem code:} catalogue entry
\subsection{Definitions and mathematical statement}
Let \(\widehat{\C}=\C\cup\{\infty\}\), and let \(f:\widehat{\C}\to\widehat{\C}\) be a nonconstant rational map of degree \(d\). The supplied statement does not explicitly assume \(d\geq 2\). An analytic Jordan curve is understood here as the image of a regular real-analytic embedding of the circle into \(\widehat{\C}\); this regularity convention requires confirmation from the source.

Consider an analytic Jordan curve \(\gamma\subset\widehat{\C}\) for which there exists an integer \(p\geq 1\) such that
\[
f^p(\gamma)=\gamma,
\qquad h=f^p|_\gamma:\gamma\to\gamma
\]
is an orientation-preserving homeomorphism with irrational Poincar\'e rotation number. More explicitly, after identifying the oriented curve with \(\R/\Z\), a lift \(H:\R\to\R\) of \(h\) satisfies \(H(t+1)=H(t)+1\), and
\[
\rho(h)=\left[\lim_{n\to\infty}\frac{H^n(t)-t}{n}\right]\in(\R/\Z)\setminus(\Q/\Z).
\]
Only irrationality is used, so reversing the chosen orientation does not affect the condition.

A Herman ring of \(f\) means a periodic doubly connected Fatou component \(U\) such that, for its period \(q\geq 1\), the restriction \(f^q|_U\) is conformally conjugate to an irrational rigid rotation of a round annulus \(\{z:r<|z|<1\}\), where \(0<r<1\). Require additionally that
\[
\gamma\not\subseteq\overline U
\quad\text{for every Herman ring }U\text{ of }f,
\]
where closure is taken in \(\widehat{\C}\). This interprets the source's exclusion as noncontainment; it does not require disjointness from every such closure. Let \(\mathcal D(f)\) denote the collection of curves satisfying these conditions, the candidates called degenerate Herman rings in the supplied problem. No separate definition of that term is supplied.

The existence question is whether there exist \(f\) and \(\gamma\in\mathcal D(f)\) such that \(\gamma\) is not a circle. The source leaves the circle convention unspecified: in a chosen affine coordinate, an ordinary circle has the form \(\{z:|z-a|=R\}\), with \(a\in\C\) and \(R>0\); a generalized circle on \(\widehat{\C}\) may also be a straight line together with \(\infty\). The intended convention must be fixed before this exclusion has a unique meaning.

The original statement also asks whether the number of degenerate Herman rings can be estimated in terms of \(d=\deg f\); this question is absent from the cleaned statement. One possible uniform interpretation is the existence of a function \(B:\Z_{\geq 1}\to\Z_{\geq 0}\) such that
\[
\#\mathcal D(f)\leq B(\deg f)
\quad\text{for every admissible rational map }f.
\]
This is a provisional interpretation only: the source does not specify whether to count individual curves, periodic cycles of curves, or only noncircular examples, nor whether the intended degree range is \(d\geq 2\). These choices remain unresolved.
\subsection{Short English statement}
Ask whether a rational map can have a noncircular analytic periodic Jordan curve on which an iterate is a homeomorphism with irrational rotation number, while the curve is contained in no Herman-ring closure. The original record also asks for an estimate of their number in terms of the map's degree. The circle convention, degree range, meaning of degenerate Herman ring, and counting convention require clarification. The supplied partial-resolution assessment concerns smooth examples and distinguishes the stronger analytic question.
\subsection{Sources}
\begin{itemize}
\item Local UnsolvedMath registry, permanent ID \(30001390\), source problem code \texttt{catalogue entry}, title \emph{Existence of Noncircular Degenerate Herman Rings}. The supplied record lists proposed year 2009 and supplies no named proposer.
\item Supplied primary citation: \emph{Mini-Workshop: The Escaping Set in Transcendental Dynamics} (2010), Oberwolfach Reports, pp.\ 2927--2964, DOI \texttt{10.4171/owr/2009/54}. \url{https://doi.org/10.4171/owr/2009/54}. The supplied record attributes the problem to this workshop report and labels the statement \texttt{corrected\_verified}, reporting comparison with the official report. That verification claim is supplied provenance, not an independent verification performed for this draft.
\item Supplied literature reference: ``Yang smooth degenerate Herman rings.'' \url{https://doi.org/10.1016/j.aim.2024.109827}.
\item Supplied literature reference: ``Eremenko 2025 invariant curves problems.'' \url{https://www.math.purdue.edu/~eremenko/dvi/invariant.pdf}.
\item Supplied literature reference: ``original OWR report.'' \url{https://doi.org/10.4171/owr/2009/54}.
\item Literature assessment dated exactly \texttt{2026-08-21}: the supplied curation labels the problem partially solved, attributes the existence of smooth degenerate Herman rings to Yang's construction, and reports that Eremenko's 2025 problem list leaves the stronger analytic invariant-curve question open. This is supplied dated curation, not independently verified. The record explicitly requests independent verification of time-sensitive status and cited claims.
\end{itemize}
\end{document}
